\documentclass[french]{book}
\usepackage{teach}
\usepackage{verbatim}
\usepackage{amsmath,amsfonts,amssymb}
\usepackage{pgfplots}
\usepackage{mwe}
\RequirePackage{graphicx}
\RequirePackage{ifpdf}
 \ifpdf
   \DeclareGraphicsRule{*}{mps}{*}{}
 \fi


  \usepackage{fancyhdr}

  \usepackage{amsmath}
  \usepackage{amssymb}
  \usepackage{amsfonts}
  \usepackage{amsthm}
  \usepackage{mathrsfs}

  \usepackage{array}
  \usepackage{multicol}
  \setlength{\multicolsep}{3pt}

  \usepackage{graphicx}
  \graphicspath{{Figures/}}
  \usepackage{pstricks,pst-plot,pst-text,pst-tree,pst-eucl}
  \usepackage[all]{xy}

  \usepackage{fancybox}
  \usepackage{color}

%  \usepackage[frenchb]{babel}
  \usepackage{hyperref}

\usepackage{subfig}
\pgfplotsset{compat=newest}
\usepackage[all]{xy}
%\usepackage{geometry}
%\geometry{top=1cm, bottom=1cm, left=2cm, right=2cm}
\usepackage[utf8]{inputenc}
\usepackage[french]{babel}
\redefinecolor{thm}{title}{bgcolor}{alizarin}
\redefinecolor{prop}{title}{bgcolor}{meatbrown}
\redefinecolor{defn}{title}{bgcolor}{olivine}
\definecolor{alizarin}{rgb}{0.82, 0.1, 0.26}
\definecolor{meatbrown}{rgb}{0.9, 0.72, 0.23}
\definecolor{olivine}{rgb}{0.6, 0.73, 0.45}
\usepackage{mathrsfs}
%%
\newcommand{\R}{\mathbb {R}}
%\newcommand{\C}{\mathds {C}}
\newcommand{\N}{\mathbb {N}}
\newcommand{\Z}{\mathbb {Z}}
\newcommand{\Q}{\mathbb {Q}}
\newcommand{\D}{\mathbb {D}}
%%
\newcommand{\se}{\geq}
\newcommand{\ie}{\leq}
\newcommand{\tm}{\times}
%\newcommand{\ell}{l}
%\newcommand{\ell'}{l'}
%%
\newcommand{\ds}{\displaystyle}
\newcommand{\limn}{\lim_{n\rightarrow +\infty}}
\newcommand{\limo}[1][x]{\ds\lim_{#1\rightarrow 0}}
\newcommand{\limite}[2][x]{\ds\lim_{#1\rightarrow #2}}
\newcommand{\limpinf}[1][x]{\ds\lim_{#1\rightarrow +\infty}}
\newcommand{\limminf}[1][x]{\ds\lim_{#1\rightarrow -\infty}}
\newcommand{\limiteg}[2][x]{\ds\lim_{#1\xrightarrow{<}#2}}
\newcommand{\limited}[2][x]{\ds\lim_{#1\xrightarrow{>}#2}}
%%
\newcommand{\vect}[1]{\overrightarrow{#1}}
\newcommand{\oij}{(\textit{O},{\overrightarrow{i}},{\overrightarrow{j}})}
%
\newcommand{\cp}[2]{%
        \begin{pmatrix}
        #1\\
        #2
        \end{pmatrix}%
        }
%
\newcommand{\calb}{\mathscr{B}}
\newcommand{\calc}{\mathscr{C}}
\newcommand{\cald}{\mathscr{D}}
\newcommand{\cale}{\mathscr{E}}
\newcommand{\calf}{\mathscr{F}}
\newcommand{\calp}{\mathscr{P}}
\newcommand{\cals}{\mathscr{S}}
\newcommand{\calt}{\mathscr{T}}
%
\newcommand{\suite}[1][u]{\left( #1_{n} \right)_{n\in\N}}
\newcommand{\suitep}[1][u]{\left( #1_{n} \right)_{n\in\N^{*}}}
\usetikzlibrary{arrows}

\newcommand{\poly}[1][a]{#1_0+#1_1x+\cdots+#1_nx^n}


\begin{document}
\tableofcontents
\chapter{Dérivation des fonctions}
\section{Généralités}
\subsection{Limite en 0}
\begin{defn}
Soit $f$ une fonction définie sur un intervalle $I$ contenant 0 et $L \in \R$.

On dit que $f(x)$ tend vers $L$ quand $x$ tend vers 0 si on peut rendre $f(x)$ aussi proche de $L$ que l'on veut pour $x$ suffisamment proche de zéro. On note : $\limo[h] f(x)=L$
\end{defn}

\begin{exemple}
$\ds \limo[x]x^{2}=0$ \hspace{1em} $\ds \limo[x]x+1=1$ \hspace{1em} $\ds \limo[x]\sqrt{x}+x-2=-2$...

Déterminer la limite en 0 de $\dfrac{x^{2}-2x}{3x}$.
\end{exemple}


Pour tout $x\neq 0$, $\dfrac{x^{2}-2x}{3x}=\dfrac{x-2}{3}$ donc $\ds \limo[x]\dfrac{x^{2}-2x}{3x}=-\dfrac{2}{3}$.

%
\subsection{Nombre dérivé}
\begin{defn}
Soit $f$ une fonction définie sur un intervalle $I$. Soit $a \in I$ et $h \neq 0$ tel que $a+h \in I$.

On appelle taux de variation de $f$ entre $a$ et $h$ le réel $\dfrac{f(a+h)-f(a)}{h}$.

On dit que $f$ est dérivable en $a$ si, lorsque $h$ tend vers 0, le taux de variation de $f$ entre $a$ et $h$ tend vers un réel $L$ autrement dit si $\ds \lim_{h\rightarrow 0}\dfrac{f(a+h)-f(a)}{h}=L$.

Ce réel $L$ est appelé nombre dérivé de $f$ en $a$ et se note $f'(a)$.
\end{defn}

\begin{exemple}
Démontrer que la fonction $f : x \longmapsto x^{2}$ est dérivable en 2 et calculer $f'(2)$.

Démontrer que la fonction $g : x \longmapsto \sqrt{x}$ n'est pas dérivable en 0.
\end{exemple}


Pour tout $h\neq 0$, $\dfrac{f(2+h)-f(2)}{h}=\dfrac{(2+h)^{2}-4}{h}=\dfrac{h^{2}+4h}{h}=h+4$

Ainsi $\ds \limo[h]\dfrac{f(2+h)-f(2)}{h}=4$.

$f$ est donc dérivable en 2 et $f'(2)=4$.

Pour tout $h\neq 0$, $\dfrac{g(0+h)-g(0)}{h}=\dfrac{\sqrt{h}}{h}=\dfrac{1}{\sqrt{h}}$.

Or, $\dfrac{1}{\sqrt{h}}$ n'a pas de limite finie lorsque $h$ tend vers 0 donc $g$ n'est pas dérivable en 0.


\subsection{Interprétation graphique}
\begin{prop}
Soit $f$ une fonction définie sur $I$ et $C_{f}$ sa représentation graphique dans un repère. Soit $a\in I$.

Si $f$ est dérivable en $a$ alors $f'(a)$ est le coefficient directeur de la tangente à $C_{f}$ au point
$A(a;f(a))$. L'équation de cette tangente est alors : $y=f'(a)(x-a)+f(a)$.
\end{prop}

\begin{minipage}{9cm}
\underline{Illustration :} soit $A(a;f(a))\in C_{f}$.\par Soit $h \neq 0$ et $M(a+h;f(a+h))\in C_{f}$.

Le quotient $\dfrac{f(a+h)-f(a)}{h}$ représente le coefficient directeur de la droite $(AM)$. Lorsque $h$ tend vers 0, $M$ se rapproche de $A$ et la droite $(AM)$ tend à se confondre avec la tangente à $C_{f}$ en $A$.
\end{minipage}
\hfill
\begin{minipage}{5.6cm}
\begin{center}
\includegraphics{Figures/CoursDerivationFigures.1}
\end{center}
\end{minipage}
%\newpage

\begin{demo}
Soit $T$ la tangente à la courbe au point $A$.
Son coefficient directeur est $f'(a)$ donc l'équation réduite de $T$ est de la forme $y=f'(a)x+b$.

$A\in T$ donc $f(a)=f'(a)a+b$ donc $b=f(a)-af'(a)$.

L'équation de $T$ est donc $y=f'(a)x+f(a)-af'(a)$ soit $y=f'(a)(x-a)+f(a)$.
\end{demo}

\subsection{Approximation affine}
\begin{prop}
Soit $f$ une fonction définie sur $I$ et dérivable en $a\in I$.
\begin{itemize}
\item Il existe une fonction $\varphi$ telle que pour tout réel $h$ avec $a+h\in I$ :
\[f(a+h)=f(a)+hf'(a)+h\varphi(h)\hspace{2em}\text{et}\hspace{2em}\limo[h]\varphi(h)=0\]
\item La fonction $h \longmapsto f(a)+hf'(a)$ est une approximation affine de $f$ pour $h$ proche de 0.
\end{itemize}
\end{prop}

\begin{demo}
Pour $h\neq 0$, on pose $\varphi(h)=\dfrac{f(a+h)-f(a)}{h}-f'(a)$.

$f$ est dérivable en $a$ donc lorsque $h$ tend vers 0, $\varphi(h)$ tend vers $f'(a)-f'(a)=0$.

De plus, $h\varphi(h)=f(a+h)-f(a)+hf'(a)$ soit $f(a+h)=f(a)+hf'(a)+h\varphi(h)$
\end{demo}

\newpage
\section{Calculs de dérivées}
\subsection{Fonction dérivée}
\begin{defn}
Soit $f$ une fonction défine sur un intervalle $I$.
\begin{itemize}
\item Si, pour tout $x$ de $I$, $f$ est dérivable en $x$, on dit que $f$ est dérivable sur $I$.
\item La fonction définie sur $I$ par $x \longmapsto f'(x)$ est appelée fonction dérivée de $f$. Cette fonction est notée $f'$.
\end{itemize}
\end{defn}

\subsection{Dérivées usuelles}
\subsubsection*{Fonctions constantes}
\begin{prop}

Si $f$ est la fonction définie sur $\R$ par :  $f(x)=k$ \\
alors $f$ est dérivable sur $\R$ et pour tout réel $x$ :  $f'(x)=0$.

\end{prop}

\begin{demo}
Pour tout réel $a$ et $h\neq 0$, 
\[\dfrac{f(a+h)-f(a)}{h}=\dfrac{k-k}{h}=0\]

$f$ est donc dérivable en $a$ et $f'(a)=0$.
\end{demo}

\subsubsection*{Fonctions affines}
\begin{prop}
Si $f$ est la fonction définie sur $\R$ par :  $f(x)=mx+p$ \\
alors $f$ est dérivable sur $\R$ et pour tout réel $x$ :  $f'(x)=m$.
\end{prop}

\begin{demo}
Pour tout réel $a$ et $h\neq 0$,
\[\dfrac{f(a+h)-f(a)}{h}=\dfrac{m(a+h)+p-ma-p}{h}=\dfrac{mh}{h}=m\]

$f$ est donc dérivable en $a$ et $f'(a)=m$
\end{demo}

\subsubsection*{Fonction carré}
\begin{prop}
Si $f$ est la fonction définie sur $\R$ par :  $f(x)=x^{2}$ \\
alors $f$ est dérivable sur $\R$ et pour tout réel $x$ :  $f'(x)=2x$.

\end{prop}

\begin{demo}
Pour tout réel $a$ et $h\neq 0$,
\[\dfrac{f(a+h)-f(a)}{h}=\dfrac{(a+h)^{2}-a^{2}}{h}=\dfrac{2ah+h^{2}}{h}=2a+h\]

De plus, $2a+h$ tend vers $2a$ lorsque $h$ tend vers 0.
$f$ est donc dérivable en $a$ et $f'(a)=2a$.
\end{demo}

\subsubsection*{Fonctions puissances}
\begin{prop}
Soit $n$ un entier tel que $n \se 1$.

Si $f$ est la fonction définie sur $\R$ par :  $f(x)=x^{n}$ \\
alors $f$ est dérivable sur $\R$ et pour tout réel $x$ :  $f'(x)=nx^{n-1}$.

\end{prop}

\begin{demo}
On rappel que $(x+h)^n=\sum^n_{k=0} C^k_n x^{n-k}h^k$, avec $ C^k_n \in \N$ le nombre de combinaison de $k$ parmi $n$ où $C^k_n=\dfrac{n!}{(n-k)!k!}$.\\

Par suite,  soit $f$ est la fonction définie sur $\R$ par :  $f(x)=x^{n}$, on a alors 
$$\dfrac{f(x+h)-f(x)}{h}=\dfrac{(x+h)^n -x^n }{h}=\dfrac{\sum^n_{k=0} C^k_n x^{n-k}h^k -x^n}{h}= \dfrac{x^n+nx^{n-1}h+\sum^n_{k=2} C^k_n x^{n-k}h^k -x^n}{h}$$
En développant les deux premiers termes de la somme on obtient alors
$$\dfrac{x^n+nx^{n-1}h+\sum^n_{k=2} C^k_n x^{n-k}h^k -x^n}{h}=\dfrac{nx^{n-1}h+\sum^n_{k=2} C^k_n x^{n-k}h^k}{h}$$
puis en divisant chaque terme par $h$ l'on a:
$$nx^{n-1}+ \dfrac{\sum^n_{k=2} C^k_n x^{n-k}h^k}{h}=nx^{n-1}+\sum^n_{k=2} C^k_n x^{n-k}h^{k-1}$$
Or ici $k\se2$ donc $k-1\se1$ ainsi le facteur $h$ apparait au moins une fois dans chacun des termes. Ainsi pour tout $k\se2$,  $\limo[h] C^k_n x^{n-k}h^{k-1} = 0$ et donc en sommant pour $k$ allant de $2$ à $n$ on obtient
$$\limo[h]  \sum^n_{k=2} C^k_n x^{n-k}h^{k-1}=0$$
Il s'ensuit que 
$$ \limo[h] \dfrac{f(x+h)-f(x)}{h}=\limo[h] \dfrac{(x+h)^n -x^n }{h}= \limo nx^{n-1}+\sum^n_{k=2} C^k_n x^{n-k}h^{k-1} = nx^{n-1} $$
\end{demo}

\subsubsection*{Fonction inverse}
\begin{prop}

Si $f$ est la fonction définie sur $]-\infty;0[~\cup~]0;+\infty[$ par :  $f(x)=\dfrac{1}{x}$ \\
alors $f$ est dérivable sur $]-\infty;0[$ et sur $]0;+\infty[$ et pour tout $x\neq 0$ :  $f'(x)=-\dfrac{1}{x^{2}}$

\end{prop}

\begin{demo}
Pour tout réel $a\neq 0$ et $h\neq 0$, \[\dfrac{f(a+h)-f(a)}{h}=\dfrac{\dfrac{1}{a+h}-\dfrac{1}{a}}{h}=\dfrac{a-(a+h)}{ha(a+h)}=\dfrac{-1}{a(a+h)}\]

Or $\dfrac{-1}{a(a+h)}$ tend vers $-\dfrac{1}{a^{2}}$ lorsque $h$ tend vers 0.

$f$ est donc dérivable en $a$ et $f'(a)=-\dfrac{1}{a^{2}}$
\end{demo}

\subsubsection*{Fonction racine carrée}
\begin{prop}

Si $f$ est la fonction définie sur $[0~;~+\infty[$ par :  $f(x)=\sqrt{x}$ \\
alors $f$ est dérivable sur $]0~;~+\infty[$ et pour tout réel $x>0$ :  $f'(x)=\dfrac{1}{2\sqrt{x}}$.

\end{prop}

\begin{demo}
Pour tout réel $a>0$ et $h\neq 0$ tel que $a+h>0$,
\begin{multline*}
\dfrac{f(a+h)-f(a)}{h}=\dfrac{\sqrt{a+h}-\sqrt{a}}{h}=\dfrac{(\sqrt{a+h}-\sqrt{a})(\sqrt{a+h}+\sqrt{a})}{h(\sqrt{a+h}+\sqrt{a})}\\
=\dfrac{a+h-a}{h(\sqrt{a+h}+\sqrt{a})}=\dfrac{1}{\sqrt{a+h}+\sqrt{a}}
\end{multline*}
Or, lorsque $h$ tend vers 0, $\dfrac{1}{\sqrt{a+h}+\sqrt{a}}$ tend vers $\dfrac{1}{2\sqrt{a}}$

$f$ est donc dérivable en $a$ et $f'(a)=\dfrac{1}{2\sqrt{a}}$
\end{demo}

\subsubsection*{Fonctions trigonométriques}
\begin{prop}
Les fonctions sinus et cosinus sont dérivables sur $\R$ et pour tout réel $x$ :
\[\sin'(x)=\cos(x)\hspace{2em}\text{et}\hspace{2em}\cos'(x)=-\sin(x)\]
\end{prop}

Résultat admis.

\subsection{Opérations sur les fonctions et dérivées}
$u$ et $v$ désignent deux fonctions dérivables sur un intervalle $I$ et $\lambda$ un réel.

\begin{prop}
La fonction $u+v$ est dérivable sur $I$ et  $(u+v)'=u'+v'$.\\
La fonction $\lambda u$ est dérivable sur $I$ et  $(\lambda u)'=\lambda u'$.
\end{prop}

\begin{demo}
Pour tout $a\in I$ et $h\neq 0$ tel que $a+h\in I$ :

\begin{align*}
\dfrac{(u+v)(a+h)-(u+v)(a)}{h}&=\dfrac{u(a+h)+v(a+h)-u(a)-v(a)}{h}\\
&=\dfrac{u(a+h)-u(a)}{h}+\dfrac{v(a+h)-v(a)}{h}
\end{align*}

dont la limite est $u'(a)+v'(a)$ lorsque $h$ tend vers 0.

\[\dfrac{(\lambda u)(a+h)-(\lambda u)(a)}{h}=\dfrac{\lambda u(a+h)-\lambda u(a)}{h}=\lambda \dfrac{u(a+h)-u(a)}{h}\]

dont la limite est $\lambda u'(a)$ lorsque $h$ tend vers 0.
\end{demo}

\begin{exemple}
Calculer la dérivée de la fonction $f$ définie sur $]0~;~+\infty[$ par $f(x)=3x^{2}+\dfrac{1}{3x}-5\sqrt{x}+2$
\end{exemple}


$f$ est dérivable sur $]0~;~+\infty[$ car elle est la somme de fonctions dérivables sur $]0~;~+\infty[$

et, pour tout $x\in]0~;~+\infty[$, $f'(x)=3\times 2x+\dfrac{1}{3}\times \left(-\dfrac{1}{x^{2}}\right)-5\times \dfrac{1}{2\sqrt{x}}$ ainsi :
\[f'(x)=6x-\dfrac{1}{3x^{2}}-\dfrac{5}{2\sqrt{x}}\]


\vspace*{2ex}
\begin{prop}
La fonction $uv$ est dérivable sur $I$ et  $(uv)'=u'v+uv'$.\\
La fonction $u^{2}$ est dérivable sur $I$ et  $(u^{2})'=2uu'$.

\end{prop}

\begin{demo}
Pour tout $a\in I$ et $h\neq 0$ tel que $a+h\in I$ :
\begin{align*}
\dfrac{(uv)(a+h)-(uv)(a)}{h}&=\dfrac{u(a+h)v(a+h)-u(a)v(a+h)+u(a)v(a+h)-u(a)v(a)}{h}\\
&=\underbrace{\dfrac{u(a+h)-u(a)}{h}}_{\text{tend vers } u'(a)}\times v(a+h)+\underbrace{\dfrac{v(a+h)-v(a)}{h}}_{\text{tend vers } v'(a)}\times u(a)
\end{align*}

dont la limite est $u'(a)v(a)+v'(a)u(a)$ lorsque $h$ tend vers 0.

Ainsi $uv$ est dérivable et $(uv)'=u'v+uv'$.
\end{demo}

\begin{exemple}
Calculer la dérivée de la fonction $f$ définie sur $\R$ par $f(x)=(x^{2}+1)(x^{5}+2)$
\end{exemple}


$f(x)=u(x)v(x)$ avec $u(x)=x^{2}+1$ et $v(x)=x^{5}+2$.

$f$ est donc dérivable sur $\R$ en tant que produit de fonctions dérivables sur $\R$ et
\[f'(x)=2x(x^{5}+2)+5x^{4}(x^{2}+1)\]


\begin{prop}
Si $v$ ne s'annule pas sur $I$

La fonction $\dfrac{1}{v}$ est dérivable sur $I$ et  $\left(\dfrac{1}{v}\right)'=-\dfrac{v'}{v^{2}}$\\
La fonction $\dfrac{u}{v}$ est dérivable sur $I$ et  $\left(\dfrac{u}{v}\right)'=\dfrac{u'v-uv'}{v^{2}}$
\end{prop}

\begin{demo}
Pour tout $a\in I$ et $h\neq 0$ tel que $a+h\in I$ :
\[\dfrac{\dfrac{1}{v(a+h)}-\dfrac{1}{v(a)}}{h}=\dfrac{v(a)-v(a+h)}{hv(a+h)v(a)}=-\dfrac{v(a+h)-v(a)}{h}\times \dfrac{1}{v(a)v(a+h)}\]

dont la limite est $-v'(a)\times \dfrac{1}{(v(a))^{2}}$ lorsque $h$ tend vers 0.

Ainsi $\dfrac{1}{v}$ est dérivable et $\left(\dfrac{1}{v}\right)'=-\dfrac{v'}{v^{2}}$

De plus \[\left(\dfrac{u}{v}\right)'=\left(u\times \dfrac{1}{v}\right)'=u'\times \dfrac{1}{v}+u\times \left(-\dfrac{v'}{v^{2}}\right)=\dfrac{u'v-uv'}{v^{2}}\]
\end{demo}

\begin{exemple}
Déterminer la dérivée de la fonction $f$ définie sur $]2~;~+\infty[$ par $f(x)=\dfrac{3}{2x-4}$ et de la fonction $g$ définie sur $\R$ par $g(x)=\dfrac{3x+5}{2x^{2}+1}$
\end{exemple}


Pour tout $x\in]2~;~+\infty[$, $f(x)=\lambda\times \dfrac{1}{u(x)}$ avec $\lambda=3$ et $u(x)=2x-4$.

$f$ est donc dérivable sur $]2~;~+\infty[$ et $f(x)=3\times \left(-\dfrac{2}{(2x-4)^{2}}\right)=\dfrac{-6}{(2x-4)^{2}}$.

Pour tout réel $x$ , $g(x)=\dfrac{u(x)}{v(x)}$ avec $u(x)=3x+5$ et $v(x)=2x^{2}+1$.

$g$ est donc dérivable sur $\R$ et \[g(x)=\dfrac{3(2x^{2}+1)-4x(3x+5)}{(2x^{2}+1)^{2}}=\dfrac{-6x^2-20x+3}{(2x^{2}+1)^{2}}\]


\begin{prop}
% Soit $f$ une fonction définie et dérivable sur un intervalle $I$. Soient $a$ et $b$ deux réels.
% 
% Soit $J$ l'intervalle formé des réels $x$ tels que $ax+b$ appartient à $I$.
% 
% La fonction $g$ définie par $g(x)=f(ax+b)$ est dérivable sur $J$ et :
% 
% \hspace*{2cm} pour tout $x\in J$, $g'(x)=af'(ax+b)$.
Soit $f$ une fonction définie et dérivable sur un intervalle $J$ et $u$ une fonction définie et dérivable sur un intervalle $I$ telle que, pour tout $x$ de $I$, $u(x)\in J$.

La fonction $f \circ u$ est dérivable et, pour tout réel $x$ de $I$ : \[(f\circ u)'(x)=u'(x)\times f'(u(x))\]
\end{prop}


\begin{demo}
Soit $f$ une fonction définie et dérivable sur un intervalle $J$ et $u$ une fonction définie et dérivable sur un intervalle $I$ telle que, pour tout $x$ de $I$, $u(x)\in J$.

$$\dfrac{f\circ u (x+h) - f\circ u (x)}{h} = \dfrac{f\circ u (x+h) - f\circ u (x)}{h} \times \underbrace{\dfrac{u(x+h)-u(x)}{u(x+h)-u(x)}}_{=1} $$
 
 Puisque le produit est commutatif on a en commutant les dénominateurs l'équation suivante. 
 
$$\dfrac{f\circ u (x+h) - f\circ u (x)}{h} = \dfrac{f\circ u (x+h) - f\circ u (x)}{u(x+h)-u(x)} \times \dfrac{u(x+h)-u(x)}{h} $$

Or la limite quand $h$ tend vers 0 de  $\dfrac{f\circ u (x+h) - f\circ u (x)}{u(x+h)-u(x)}$ existe car $f$ est dérivable sur l'intervalle $J$ et  $u(x)\in J$ et vaut $f'(u(x))$ par définition.\\

De même, la limite quand $h$ tend vers 0 de  $\dfrac{u(x+h)-u(x)}{h}$ existe car $u$ est dérivable sur l'intervalle $I$ et $x\in I$ et vaut $u'(x)$ par définiton.\\

Par suite la limite quand $h$ tend vers 0 de  $\dfrac{f\circ u (x+h) - f\circ u (x)}{u(x+h)-u(x)} \times \dfrac{u(x+h)-u(x)}{h}$ existe et vaut $f'(u(x)) \times u'(x) = u'(x)\times f'(u(x))$. \\ 

Autrement dit, la limite quand $h$ tend vers 0 de  $\dfrac{f\circ u (x+h) - f\circ u (x)}{h}$ existe et est égale à $u'(x)\times f'(u(x))$. \\

Ce qui équivaut encore à dire,

 $$ (f\circ u)'(x)=u'(x)\times f'(u(x))$$


\end{demo}



\begin{exemple}
Calculer la dérivée de la fonction $g$ définie sur $\R$ par $g(x)=\cos\left(2x+\dfrac{\pi}{3}\right)$.
\end{exemple}


Pour tout réel $x$, $g(x)=f\circ u (x)$ avec $u(x)=2x+\dfrac{\pi}{3}$ et $f(x)=\cos(x)$.

$g$ est donc dérivable sur $\R$ et $g'(x)=-2\sin\left(2x+\dfrac{\pi}{3}\right)$


\begin{defn}
Soit $f$ une fonction définie et dérivable sur un intervalle $J$, on dit que $f$ est deux fois dérivable sur $J$ si la dérivée de $f'$ existe sur $J$. 
\end{defn}


\section{Applications de la dérivation aux variations de fonctions}
\subsection{Dérivée et variations}
\begin{prop}
Soit $f$ une fonction dérivable sur un intervalle $I$.
\begin{itemize}
\item $f$ est croissante sur $I$ si et seulement si pour tout $x$ de $I$, $f'(x)\se 0$.
\item $f$ est constante sur $I$ si et seulement si pour tout $x$ de $I$, $f'(x)=0$.
\item $f$ est décroissante sur $I$ si et seulement si pour tout $x$ de $I$, $f'(x)\ie 0$.
\end{itemize}
\end{prop}

Propriété admise.

\begin{rem}
on utilise souvent les résultats suivants.
\begin{itemize}
\item Si, pour tout $x$ de $I$, $f'(x)>0$ alors $f$ est strictement croissante sur $I$.
\item Si, pour tout $x$ de $I$, $f'(x)<0$ alors $f$ est strictement décroissante sur $I$.
\end{itemize}
\end{rem}

\begin{exemple}
Étudier les variations de la fonction $f$ définie sur $\R$ par $f(x)=x^{3}+2x$.
\end{exemple}


$f$ est dérivable sur $\R$ et pour tout $x\in \R$, $f'(x)=3x^{2}+2>0$.

La fonction $f$ est donc strictement croissante sur $\R$.


\subsection{Extremum local}
\begin{prop}
Soit $f$ une fonction définie sur un intervalle $I$ et $x_{0}\in I$.

On dit que $f(x_{0})$ est un maximum local (respectivement minimum local) de $f$ si l'on peut trouver un intervalle ouvert $J$ inclus dans $I$ et contenant $x_{0}$ tel que pour tout $x\in J$, $f(x)\ie f(x_{0})$ (respectivement $f(x)\se f(x_{0})$).
\end{prop}

\begin{minipage}{7.5cm}
\begin{exemple}
Une fonction est représentée ci-contre.

Son minimum est $-2$

Son maximum est 2.

$-1$ et $-2$ sont des minimums locaux

1 et 2 sont des maximums locaux.
\end{exemple}
\end{minipage}
\hfill
\begin{minipage}{6.5cm}
\begin{center}
\includegraphics{Figures/CoursDerivationFigures.2}
\end{center}
\end{minipage}

\begin{prop}
Soit $f$ une fonction dérivable sur un intervalle $I$ et $x_{0}\in I$.

Si $f(x_{0})$ est un extremum local de $f$ alors $f'(x_{0})=0$.
\end{prop}


\begin{minipage}{7.5cm}
\begin{rem}
 La réciproque de cette propriété est fausse, voir le contre-exemple sur la figure çi contre.
\end{rem}
\end{minipage}
\hfill
\begin{minipage}{6.5cm}
\begin{center}
\includegraphics{Figures/coursderivation3.1}
\end{center}
\end{minipage}

\begin{prop}
Soit $f$ une fonction dérivable sur un intervalle $I$ et $x_{0}\in I$.

Si $f'$ s'annule en changeant de signe en $x_{0}$ alors $f(x_{0})$ est un extremum local de $f$.
\end{prop}


\begin{defn}
Soit $f$ une fonction définie et dérivable sur un intervalle $J$, on dit que $f$ est deux fois dérivable sur $J$ si la dérivée de $f'$ existe sur $J$. 
\end{defn}

\newpage

\section{Applications de la dérivation aux fonctions convexes}
\subsection{Convexité $ \leftrightharpoons$ Concavité}

\begin{defn}
\begin{itemize}
\item[$\bullet$] Une fonction $f$, définie, dérivable (donc continue) sur un intervalle $I$ est convexe sur $I$ si sa courbe représentative est entièrement située au-dessus de chacune de ses tangentes.\\

\item[$\bullet$] Une fonction $f$, définie, dérivable (donc continue) sur un intervalle $I$ est concave sur $I$ si sa courbe représentative est entièrement située en-dessous de chacune de ses tangentes.
\end{itemize}
\end{defn}

\begin{prop}
Soit $f$ une fonction définie et deux fois dérivable sur un intervalle $J$ alors les trois propositions sont équivalentes: 

\begin{itemize}
\item[$\bullet$] La fonction $f$ est convexe.
\item[$\bullet$]  La courbe représentative $C_f$ est entièrement située au-dessus de chacune de ses tangentes.
\item[$\bullet$]  La fonction dérivé $f'$ est croissante sur $J$.
\item[$\bullet$] La dérivée seconde $f''$ est positive.
\end{itemize}

On peut trouver de manière analogue des propriétés sur la concavité.

\begin{itemize}
\item[$\bullet$] La fonction $f$ est concave.
\item[$\bullet$]  La courbe représentative $C_f$ est entièrement située en-dessous de chacune de ses tangentes.
\item[$\bullet$]  La fonction dérivé $f'$ est décroissante sur $J$.
\item[$\bullet$] La dérivée seconde $f''$ est négative.
\end{itemize}

\end{prop}

\subsection{Point d'inflexion}

Certaines fonctions sont convexes sur un intervalle puis concaves sur un autre intervalle \emph{(ou inversement)}. On cherche à savoir à partir de quels $x$ il y a un changement entre convexité et concavité. Pour cela on remarque qu'une fonction continue ne peut changer entre convexité et concavité que quand la courbe représentative coupe la tangente. \\

\emph{Effectivement, supposons que notre fonction soit convexe et qu'elle devienne concave à partir d'un certain $x$, alors forcément la courbe représentative est au dessus de ses tangentes avant $x$, si elle devient concave alors forcément la courbe représentative est en desous de ses tangentes  après $x$, donc elle coupe nécéssaire la tangente en $x$.}

\begin{defn}
Un point d’inflexion est un point où la courbe représentative d’une fonction coupe sa tangente (et donc change de convexité).
\end{defn}

\begin{prop}[Lien avec la dérivé seconde]
Soit $f$ une fonction définie sur $I$, deux fois dérivable sur $I$ et soit $a \in I$.
Si $f''$ s’annule pour $a$ et change de signe, alors la courbe graphique de $f$ admet un point d’inflexion de coordonnées $(a ; f(a))$.
\end{prop}

\newpage
\section{Applications de la convexité à l'analyse}
\subsection{Inégalité de Jensen}


\begin{thm}[Inégalité de Jensen]
Soit $f$ une fonction convexe sur un intervalle $I$. Pour tout $(x_1,x_2, \cdots , x_n) \in I^n $ et $(\lambda_1,\lambda_2, \cdots , \lambda_n) \in ]0;1]^n $ tel que $\sum_{i=1}^n \lambda_i =1$, on a

$$f(\sum_{i=1}^{n} \lambda_i x_i) \ie \sum_{i=1}^n \lambda_i f(x_i)$$
\end{thm}

\begin{demo}
On procède par récurrence sur $\mathbb{N}*$, soient I un intervalle, et $f : I \longrightarrow \mathbb{R}$ 
convexe, $(x_1,x_2, \cdots , x_n) \in I^n $ et $(\lambda_1,\lambda_2, \cdots , \lambda_n) \in ]0;1]^n $ tel que $\sum_{i=1}^n \lambda_i =1$.\\

\textbf{Initialisation :} Si $n=1$ on a trivialement $f(\lambda x_1)\ie \lambda_1f(x_1)$ car ici $\lambda_1=1$.\\

\textbf{Hérédité : }On suppose ici que le résultat est vrai au rang $n \in \mathbb{N}^*$ fixé, montrons que ceci entraine l’inégalité pour le rang $n + 1$. On pose $T_n = \sum_{i=1}^n \lambda_i x_i $ et $S_n=\sum_{i=1}^n \lambda_i$  tels que $\dfrac{T_n}{S_n} \in I$ et $S_n + \lambda_{n+1} =1$.\\

On sait que 
$$f(\sum_{i=1}^{n+1} \lambda_i x_i =f( \sum_{i=1}^{n} \lambda_i x_i + \lambda_{n+1} x_{n+1} ) = f(T_n+ \lambda_{n+1} x_{n+1} )$$

Par convexité de $f$ on  a
$$f(S_n\dfrac{Tn}{Sn} + \lambda_{n+1} x_{n+1} ) \ie S_n f( \dfrac{T_n}{S_n}) + \lambda_{n+1} f(x_{n+1})$$ 
par hypothèse de récurrence on a alors,

$$ S_n f( \dfrac{T_n}{S_n}) + \lambda_{n+1} f(x_{n+1}) \ie S_n \sum_{i=1}^n \dfrac{\lambda_i x_i}{S_n} +  \lambda_{n+1} f(x_{n+1})$$

d'où l'inégalité suivante, 

$$f(\sum_{i=1}^{n} \lambda_i x_i) \ie \sum_{i=1}^n \lambda_i f(x_i)$$

\end{demo}

Ce théorème peut donner un résultat aussi sur les fonctions concaves rapidement puisque si $f$ est concave alors $-f$ est convexe. 

Ainsi l'on obtient un corrolaire de ce théorème.

\begin{cor}[Inégalité de Jensen pour les fonctions concaves]
Soit $f$ une fonction concave sur un intervalle $I$. Pour tout $(x_1,x_2, \cdots , x_n) \in I^n $ et $(\lambda_1,\lambda_2, \cdots , \lambda_n) \in ]0;1]^n $ tel que $\sum_{i=1}^n \lambda_i =1$, on a

$$f(\sum_{i=1}^{n} \lambda_i x_i) \se \sum_{i=1}^n \lambda_i f(x_i)$$
\end{cor}

\newpage

\subsection{Inégalité arithmético-géométrique }

Une inégalité bien connu en mathématique est l'inégalité arithmético-géométrique. Le but étant de comparer ces deux quantités pour des réels strictements positifs. Pour rappel, la moyenne arithmétique est pour des réels $x_1,x_2, \cdots , x_n$ est $\dfrac{x_1+x_2+\cdots+x_n}{n}$, et la moyenne géométriques est $\sqrt[n]{x_1x_2\cdots x_n}$.

\begin{thm}[Inégalité arithmético-géométrique]
Soit  $x_1,x_2, \cdots , x_n$ des réels strictements positifs alors 
$$\sqrt[n]{x_1x_2\cdots x_n} \ie \dfrac{x_1+x_2+\cdots+x_n}{n} $$
\end{thm}

\begin{demo}
On sait que $ln$ est une fonction concave, \emph{$ln''=(\dfrac{1}{x})'=\dfrac{-1}{x^2}$}  Posons $\lambda_1=\lambda_2=\cdots=\dfrac{1}{n}$ tel que $\sum_{i=1}^{n} \lambda_i =1$. En appliquant notre corrolaire on a donc 

$$\dfrac{1}{n} ln(x_1) + \dfrac{1}{n} ln(x_2) + \cdots +\dfrac{1}{n} ln(x_n) \ie ln( \dfrac{1}{n}x_1 +\dfrac{1}{n}x_2 + \cdots + \dfrac{1}{n}x_n )$$

Or d'après les propriétés élémentaires sur le logarithme cette inégalité est équivalente à

$$ln(\sqrt[n]{x_1x_2\cdots x_n}) \ie ln( \dfrac{x_1+x_2+\cdots+x_n}{n} )$$

L'exponentielle étant une fonction \emph{strictement} croissante, on peut l'appliquer à cette inégalité en conservant l'ordre et donc 

$$\sqrt[n]{x_1x_2\cdots x_n} \ie \dfrac{x_1+x_2+\cdots+x_n}{n} $$

qui est l'inégalité recherchée.
\end{demo}

\end{document}
